\documentclass[11pt]{article}
\usepackage[margin=1in]{geometry}
\usepackage{amsmath,amssymb,amsthm}
\usepackage[T1]{fontenc}
\usepackage{lmodern}
\usepackage{microtype}
\usepackage[colorlinks=true,allcolors=blue]{hyperref}
\hypersetup{pdftitle={Metric rigidity at the boundary of cubic compactifications},pdfauthor={Alec Kriebel},pdfsubject={Sharp Cartier-index rigidity and an attributed cubic moduli consequence}}
\setlength{\emergencystretch}{2em}
\newtheorem{proposition}{Proposition}
\newtheorem{corollary}[proposition]{Corollary}
\newtheorem{theorem}[proposition]{Theorem}
\newcommand{\nv}{\widehat{\operatorname{vol}}}
\newcommand{\C}{\mathbb C}
\newcommand{\Q}{\mathbb Q}
\newcommand{\PP}{\mathbb P}
\title{Metric rigidity at the boundary\\of cubic compactifications}
\author{Alec Kriebel\\\small ORCID: \href{https://orcid.org/0009-0001-9320-500X}{0009-0001-9320-500X}}
\date{October 6, 2026}
\begin{document}
\maketitle
\begin{abstract}
We deduce a sharp Cartier-index criterion for complex-structure rigidity
on singular K\"ahler--Einstein Fano varieties. If $-K_X=rL$ with $L$ ample
Cartier and $r>\dim(X)/2+1$, then the tangent sheaf is stable and every
parallel orthogonal complex structure on the regular locus equals the
given one or its negative. The proof combines the decomposition theorem
of Druel--Guenancia--P\u aun with a Hilbert-polynomial index bound and local
positive-Ricci curvature. Applied to cubic $n$-fold limits, $n\geq5$, it
identifies the ordinary metric Gromov--Hausdorff compactification with
cubic GIT modulo conjugation, including the singular boundary. We also
record the complex-preserving comparison as an attributed consequence
of the unrestricted ordinary-double-point gap, Li--Liu's global Reeb
minimum, and Spotti--Sun's established transfer. The complex algebraic
comparison is already publicly stated in stronger form by
Kong--Shen--Zhao--Zhao; no first-solution claim is made.
\end{abstract}

\paragraph{Status and provenance.}
The upstream manuscript \cite{OAI} is pinned at repository commit
\texttt{adc7f1241b42} (the full hash is in the source manifest). Its theorem is a
principal external input, not a theorem reproved in this note. Separate
scoped adversarial audits inspected its geometric and semigroup arguments
and its sole fourfold companion. These are evidence, not formal
certification. No Lean verification of the gap or these consequences is
claimed. AI tools were used extensively in research, drafting, and
verification. This note has not undergone conventional human peer review.
The additional assertion concerns singular metric rigidity and the fibers
of the forgetful map. The splitting principle, the conjugation convention,
and the complex cubic comparison are credited to their original sources.

\section{Exact target and dependency diagram}
Fix $n\geq5$ and let
\[
 G_n=\PP\bigl(H^0(\PP^{n+1},\mathcal O(3))\bigr)//\operatorname{SL}(n+2).
\]
The complex GH moduli object retains complex structure: its points are
biholomorphic isometry classes of K\"ahler--Einstein cubic $n$-folds,
with $\operatorname{Ric}(\omega)=\omega$, and its compactification is the
complex/polarized GH compactification used in \cite{SS}. Bare metric-space
isometry forgets complex structure. Write
$\overline{\mathcal M}_{n}^{\rm GH,metric}$ for the closure of the
underlying compact metric spaces of smooth KE cubics, with ordinary
metric GH topology. Write $\sigma$ for coefficientwise conjugation on
$G_n(\C)$. We prove that conjugation orbits are exactly the fibers of
the forgetful map on the entire boundary.

The good moduli space on the K side has closed complex points representing
K-polystable $\Q$-Fano varieties that admit a $\Q$-Gorenstein deformation
to a smooth cubic. Equivalently one considers the corresponding closure
of the smooth cubic locus in the smoothable K-moduli space, with its
complex analytic topology. Equal dimension and volume alone do not select
that smoothing locus. The established smoothable KE/GH correspondence is
the relevant identification with these closed points \cite{LWX}.

The dependency diagram for the proposed deduction is
\[
\begin{gathered}
\nv(x,X)\leq2(k-1)^k\quad(2\leq k\leq n)\quad\text{\cite{OAI}}\\
\Downarrow\quad\text{global Reeb minimum and density identity \cite{LL}}\\
\Theta(C)\leq2(1-1/k)^k\quad\text{for isolated singular metric cones}\\
\Downarrow\quad\text{iteration and Cartier-root theorem \cite{SS}}\\
-K_Z=(n-1)L_Z,\quad L_Z\text{ Cartier},\quad L_Z^n=3\\
\Downarrow\quad\text{Fujita classification; CM positivity; continuity \cite{Fujita,SS}}\\
\overline{\mathcal M}_{n}^{\rm GH,\C}\ \cong_{\rm homeo}\ G_n(\C).
\end{gathered}
\]
Only the upper gap bound is needed; its equality classification is not
used. No scheme/stack isomorphism, equivalence of moduli functors, or
assertion about every nonclosed semistable orbit is deduced here.

\begin{theorem}[Sharp Cartier-index rigidity]\label{thm:index}
Let $X$ be an irreducible normal projective klt Fano of dimension $m\geq1$
carrying a weak KE current with bounded potentials, smooth on $X_{\rm reg}$.
Suppose an ample Cartier line bundle $L$ and a positive integer $r$
satisfy the actual line-bundle identity
$\omega_X^{[-1]}\simeq L^{\otimes r}$ and $r>m/2+1$.
Then $T_X$, and the tangent sheaf of every finite quasi-\'etale cover,
is slope stable for the anticanonical polarization. No such cover is
a product of two positive-dimensional normal projective varieties.
Every parallel orthogonal complex structure on $(X_{\rm reg},g)$
is $J$ or $-J$. The strict numerical threshold is sharp in this class.
\end{theorem}

\begin{theorem}[Cubic metric boundary]\label{thm:metric}
For $n\geq5$, using the upper gap theorem of \cite{OAI}, there are natural
homeomorphisms
\[
\overline{\mathcal M}_{n}^{\rm GH,\C}\simeq G_n(\C),
\qquad
\overline{\mathcal M}_{n}^{\rm GH,metric}
\simeq G_n(\C)/\langle\sigma\rangle.
\]
The first space has the closed-point interpretation in the cubic
smoothable K-polystable coarse space described above. Every metric
isometry between its KE boundary varieties is holomorphic or
antiholomorphic; in the latter case it is a holomorphic isomorphism
to the conjugate variety. In either formulation the cubic hyperplane
Cartier line bundle is preserved.
\end{theorem}

The first comparison uses the displayed established transfer. The
additional argument for the second comparison is
\[
\begin{gathered}
\text{singular KE stable-factor cover \cite{DGP}}
\ \Longrightarrow\
\dim(Y_i)\geq r-1\ \Longrightarrow\ T_X\text{ stable}\\
\Longrightarrow\ I=\pm J
\ \Longrightarrow\
\text{all metric fibers are conjugation orbits}\\
\Longrightarrow\
\text{compact--Hausdorff quotient homeomorphism}.
\end{gathered}
\]
The inequality is about a genuine Cartier root, not an arbitrary
$\Q$-Cartier divisor. The metric isometry statement uses the regular-locus
theorem for the actual GH limits; no claim about arbitrary metric
completions of nonsmoothable weak KE varieties is needed.

\section{Cartier index and parallel complex structures}
\begin{proof}[Proof of Theorem~\ref{thm:index}]
Druel--Guenancia--P\u aun \cite[Theorem A]{DGP} supply a finite quasi-\'etale
cover $f:Y\to X$ and an isometric algebraic product of positive-dimensional
weak KE klt Fanos $Y_i$ with stable tangent sheaves. Their theorem treats
the incomplete regular locus; we do not apply complete de Rham splitting
there. Put $M=f^*L$. Quasi-\'etaleness gives
$\omega_Y^{[-1]}\simeq M^{\otimes r}$.

For each factor choose a regular point $b$ of the complementary product.
The restriction $M_i=M|_{Y_i\times\{b\}}$ is ample Cartier. On
$Y_{i,\rm reg}\times\{b\}$ the smooth product canonical formula identifies
$\omega_Y$ with $\omega_{Y_i}$ tensored with the fixed one-dimensional
space $\det(T_b^*B)$. Normal reflexive extension therefore gives
$\omega_{Y_i}^{[-1]}\simeq M_i^{\otimes r}$ globally. This also works for
the factors of any algebraic product cover: quasi-\'etale pullback is klt,
and adjoining smooth coordinates tests klt on a factor.

Let $Z$ be such a factor, of dimension $d$, with $-K_Z=rA$ and $A$
ample Cartier. For $j=1,\ldots,r-1$, klt Kawamata--Viehweg vanishing
applies to $-jA$, since $-jA-K_Z=(r-j)A$ is ample
\cite[Theorem 2.22]{KV}.
Negative ampleness gives $H^0(Z,-jA)=0$ as well. Hence the nonzero
degree-$d$ polynomial $P_Z(t)=\chi(Z,A^{\otimes t})$ has the $r-1$
distinct roots $-1,\ldots,-(r-1)$, so $d\geq r-1$.
Euler characteristic is polynomial at every integer, including negative
integers \cite[Lemma 33.45.1]{Stacks}. Two factors would imply
$m\geq2(r-1)$, contradicting the strict index bound. The DGP cover
therefore has one factor and $T_Y$ is stable.

Stability descends without a $\Q$-factoriality assumption. If a proper
saturated $\mathcal F\subset T_X$ had slope at least $\mu(T_X)$, pull it
back on the finite \'etale big open set and extend reflexively and
saturate to $\mathcal G\subset T_Y$. Both complements have codimension at
least two. Choose a general complete-intersection curve of large
anticanonical degree in that big open set. On its \'etale inverse image,
ordinary bundle degrees multiply by $\deg f$, giving
\[
\mu_{f^*(-K_X)}(\mathcal G)
=\deg(f)\mu_{-K_X}(\mathcal F),\qquad
\mu_{-K_Y}(T_Y)=\deg(f)\mu_{-K_X}(T_X).
\]
This contradicts stability upstairs. The same hypotheses hold on every
finite quasi-\'etale cover, proving its tangent stability as well.

Let $I$ be another parallel orthogonal complex structure on $X_{\rm reg}$.
The real two-form $\eta(u,v)=g(Iu,v)$ is parallel. Its $(2,0)$ part
relative to $J$ is parallel. On a KE manifold with
$\operatorname{Ric}(g)=\lambda g$, $\lambda>0$, contraction of Chern
curvature on $\Lambda^pT^{*(1,0)}$ is $-p\lambda\,\operatorname{Id}$,
up to the common curvature sign convention. A parallel section has zero
curvature action, so this contraction annihilates it. Thus no nonzero
parallel $(2,0)$ part exists. This is a local identity, requiring no
completeness or integration.

Consequently $\eta$ is of type $(1,1)$, equivalently
$-JIJ=I$, so $IJ=JI$. The parallel endomorphism $I$ of $T^{1,0}_{X_{\rm reg}}$
is holomorphic. The reflexive sheaf $\mathcal Hom(T_X,T_X)$ extends it
across the codimension-two complement; GAGA makes it algebraic
\cite[Theorem 2]{GAGA}.
A stable sheaf is simple: the kernel/image slope inequalities exclude
a nonzero endomorphism of intermediate generic rank. Characteristic
polynomial coefficients extend to global regular functions on normal
projective $X$, hence are constant. Subtracting a complex eigenvalue
then gives an endomorphism of deficient generic rank, which must be zero.
Thus $I=c\,\operatorname{Id}$. The equation $I^2=-1$ gives
$c=\pm i$, or $I=\pm J$ as a real endomorphism.

For sharpness take $X=\PP^{r-1}\times\PP^{r-1}$, $r\geq2$, and
$L=\mathcal O(1,1)$. Then $m=2r-2$ and $-K_X=rL$. The product of equally
normalized Fubini--Study metrics is KE. Conjugation on only one factor
is an isometry whose pullback complex structure is $(-J_1,J_2)$,
neither $J$ nor $-J$. This falsifies including equality in the uniform
criterion; it is not a cubic counterexample.
\end{proof}

\section{The infimum-to-metric implication}
Write $\nv(x,X)=\inf_v A_X(v)^k\operatorname{vol}(v)$ for a
boundary-zero complex algebraic klt germ of dimension $k$. Ordinary volume
is normalized by $k!$ times the leading valuation-ideal colength.

\begin{proposition}
Assume the algebraic gap bound in dimension $k\geq2$. Then an isolated
singular Ricci-flat K\"ahler cone with compact smooth Sasaki--Einstein
link has density at most $b_k=2(1-1/k)^k$.
\end{proposition}
\begin{proof}
Use the cone metric $dr^2+r^2g_Y$, with
\[
 \Theta(C(Y))=\frac{\operatorname{Vol}(Y)}{\operatorname{Vol}(S^{2k-1})}.
\]
The Sasakian cone has a normal affine algebraic completion; a
quasi-regular CR-preserving approximation constructs it as an orbifold
line-bundle cone \cite[\S3]{vC}. It is $\Q$-Gorenstein. The positive
pluricanonical Reeb weight gives klt, by
\cite[Lemma 6.1]{CS}. Rationality and $\Q$-Gorensteinness alone are not
being substituted for this klt criterion. The cone germ has zero boundary.

There is a convenient way to respect the genuine top-form hypothesis in
the setup of Li--Liu \S6. If $\pi_1(Y)$ is nontrivial, Myers' theorem
gives a finite universal covering of degree $d\geq2$. Bishop's inequality
on that cover yields $\Theta(C(Y))\leq1/d\leq1/2\leq b_k$.
If $Y$ is simply connected, the Ricci-flat punctured cone has a global
parallel holomorphic top form. Its Euler weight is $k$; it extends
reflexively across the vertex and is equivariant for the cone torus.
This satisfies the canonical-form setup of Li--Liu.

Li--Liu's Theorem 1.9 (=6.2) states that the metric Reeb valuation
$v_\xi$ globally minimizes over \emph{all} real valuations centered at
the vertex. The theorem includes irregular Reeb fields. Their Lemma 6.8
and limiting proof use quasi-regular structures converging smoothly;
these approximants are not incorrectly assumed to remain Einstein.
Lemma 6.3 and the same limiting argument give
\[
 A_X(v_\xi)=k,\qquad
 \nv_X(v_\xi)=k^k\Theta(C(Y)).
\]
The ratio to the round sphere cancels the fixed conversion between contact
volume and Riemannian link volume. Both sides are continuous in the
quasi-regular approximation. Thus, in the simply connected case,
\[
 k^k\Theta(C(Y))=\nv(o,X)\leq2(k-1)^k.
\]
An upper bound on the infimum alone would not imply this inequality for
an arbitrary Reeb valuation. The global minimum is essential.
\end{proof}

Flat $\C^k$ has density $1$ and normalized volume $k^k$; it is excluded
from the singular assertion. The ODP degree valuation has discrepancy
$k-1$ and volume $2$, consistent with the stated normalization. No direct
application of Li--Liu to a cone with singular link is made.

\section{From density to polarized cubic limits}
Spotti--Sun's constant $A'(n)$ is the supremum of the densities of
products of a flat factor with an isolated singular cone having smooth
link. Their iterated-cone argument treats other singular links
\cite[Lemma 5.3 and Theorem 5.2]{SS}. The required transverse dimensions
are exactly $k=2,\ldots,n$. A normal klt algebraic curve is smooth, and
the relevant KE singular strata have complex codimension at least two.

The sequence $b_k$ increases. For real $t>1$,
\[
 \frac{d}{dt}\log b_t
 =\log(1-1/t)+\frac1{t-1}
 =u-\log(1+u)>0,\qquad u=\frac1{t-1}.
\]
Products with Euclidean space do not change density, so the preceding
proposition gives $A'(n)\leq b_n$.

For a smooth cubic, adjunction gives
\[
 -K_X=(n-1)H,\qquad H^n=3,\qquad V=(-K_X)^n=3(n-1)^n.
\]
Spotti--Sun Theorem 5.2 requires
\[
 V>\tfrac12 A'(n)(n+1)^n.
\]
This follows strictly from
\[
 \tfrac12 b_n(n+1)^n
 =(n-1)^n(1+1/n)^n<e(n-1)^n<3(n-1)^n.
\]
Its conclusion supplies both Gorenstein canonical singularities and the
Cartier root $L_Z$ with $-K_Z=(n-1)L_Z$. Preserving canonical singularities
alone would not preserve the required polarization. The theorem's proof
uses the polarized smoothability of the limit, the tangent-cone Cartier
index comparison, and quotient-smoothing/rigidity arguments. These are
inherited hypotheses and conclusions, not consequences of density alone
for arbitrary $\Q$-Fanos.

Fujita's Gorenstein del Pezzo classification applies to $(Z,L_Z)$
\cite{Fujita}. Its intermediate-cohomology hypothesis holds for all
integer twists: Kawamata--Viehweg vanishing applies to $tL_Z$ for
$t\geq2-n$, since $tL_Z-K_Z=(t+n-1)L_Z$ is ample. Serre duality applies
the same argument at twist $-t-(n-1)$ for $t\leq-1$. These ranges cover
all integers for $n\geq3$, and klt makes $Z$ Cohen--Macaulay.
The degree-three case has $h^0(Z,L_Z)=n+2$ and $L_Z$ very ample, so it
embeds $Z$ as a cubic hypersurface in $\PP^{n+1}$. This uses the canonical
Gorenstein version, without imposing an extra terminality assumption.

\section{Moduli continuity and the precise consequence}
Here is an explicit positivity check for the relevant CM polarization.
Put $S=\PP(H^0(\PP^{n+1},\mathcal O(3)))$, let $h$ and $s$ denote
the hyperplane classes on $\PP^{n+1}$ and $S$, and put $r=n-1$.
The universal cubic has class $3h+s$, and its relative anticanonical
class is $rh-s$. In the standard Fano normalization,
\[
\begin{split}
 c_1(\lambda_{\rm CM})
 &=-\pi_*\bigl((rh-s)^{n+1}(3h+s)\bigr)\\
 &=r^n\bigl(3(n+1)-r\bigr)s
 =2(n+2)(n-1)^n s.
\end{split}
\]
The expression is computed on the product; the universal-cubic factor
implements pushforward from the hypersurface. Any positive conventional
rescaling preserves the sign. The coefficient is strictly positive,
so the CM polarization agrees in sign with the standard GIT polarization.

Spotti--Sun \S5.2 applies their moduli continuity method from \S4.2.
The initial Fermat cubic is KE; the smooth cubic moduli locus is connected.
KE existence is open there, and the cubic-limit and CM facts make it
closed. The map from the complex GH compactification to GIT is continuous;
uniqueness of the KE metric on a fixed Fano variety gives injectivity.
Its compact image contains the dense smooth GIT locus, giving
surjectivity. Compactness and the Hausdorff analytic GIT topology give a
homeomorphism. This is precisely the established transfer
\cite[Theorem 1.3(2)]{SS}, rather than a new continuity theorem.

\begin{corollary}[Audited conditional consequence]
If the unrestricted algebraic normalized-volume gap holds in every
dimension $2\leq k\leq n$, then for $n\geq5$ the complex GH
compactification of smooth KE cubic $n$-folds is naturally homeomorphic
to $G_n(\C)$. Its closed points also identify with the corresponding
smoothable K-polystable Fano limits in the cubic smoothing locus.
\end{corollary}

Taking the upper-bound theorem of \cite{OAI} as the external input removes
the dimension-by-dimension metric-gap hypothesis in \cite[Theorem
1.3(2)]{SS}. This is the inherited complex-moduli deduction. The
topological result alone does not upgrade this deduction to schemes,
stacks, or nonclosed semistable points.

\section{Isometries on the singular cubic boundary}
\begin{proof}[Proof of Theorem~\ref{thm:metric}]
The complex comparison is the preceding corollary with \cite{OAI}.
Its KE boundary varieties are normal klt cubic hypersurfaces, and
adjunction gives $-K_X=(n-1)H_X$ with $H_X$ ample Cartier.
Since $n-1>n/2+1$ for $n\geq5$, Theorem~\ref{thm:index} applies.

Donaldson--Sun \cite[Proposition 2.14]{DSII} identify the metric singular
set and the analytic singular set in this GH-limit setting. A metric
isometry $F:X\to Z$ therefore maps $X_{\rm reg}$ onto $Z_{\rm reg}$.
Myers--Steenrod makes it a smooth Riemannian isometry locally there;
this local assertion requires no completeness of the regular locus.
Thus $I=F^*J_Z$ is parallel and orthogonal. Theorem~\ref{thm:index}
gives $I=J_X$ or $-J_X$ with a single global sign. Hence $F$ is
holomorphic on the regular locus, or holomorphic to $\overline Z$.

To extend it, choose an analytic chart about $F(x)$ and a neighborhood
of $x$ whose continuous image lies in that chart. Each target coordinate
composed with $F$ is holomorphic on the regular locus and locally
bounded. Normal Riemann extension makes these functions holomorphic
on the whole neighborhood. They equal the continuous original
coordinates by density and satisfy the target equations. Repeat for
$F^{-1}$. This gives a global holomorphic isomorphism, algebraic by
GAGA \cite[Proposition 15]{GAGA}, with target $Z$ or $\overline Z$.
The metric and analytic
topologies agree in the GH algebraization; the local extension
principle also appears in \cite[Proposition 2.4]{DSII}.

Grothendieck's Picard Lefschetz theorem
\cite[Expos\'e XII, Corollaire 3.7]{SGA2} applies to global projective
complete intersections of dimension at least three, including singular
ones. It gives $\operatorname{Pic}(X)=\mathbb Z[H_X]$.
Canonical preservation implies
$(n-1)(F^*H_Z-H_X)=0$, taking $\overline Z$ in the antiholomorphic
case. The Picard group is torsion free, so $F^*H_Z=H_X$.
The complete hyperplane system is the cubic embedding, since the
hypersurface sequence gives $h^0(X,H_X)=n+2$. Accordingly these
isomorphisms give precisely the corresponding projective cubic
GIT points and their conjugates.

The forgetful map from the complex GH compactification to ordinary
metric GH space is continuous. Any metric limit of smooth KE cubics
has a polarized subsequential lift by polarized compactness, so its
image is exactly $\overline{\mathcal M}_{n}^{\rm GH,metric}$.
The preceding isometry classification makes its fibers $\sigma$-orbits.
Conversely complex conjugation changes both $J$ and the real K\"ahler
two-form by sign, preserving the Riemannian metric. The induced map
from $G_n(\C)/\langle\sigma\rangle$ is therefore a continuous
bijection from a compact space to a Hausdorff space, hence a
homeomorphism.
\end{proof}

No effective GH-distance estimate or modulus of continuity is asserted.
The sharpness example in Theorem~\ref{thm:index} concerns general Fanos;
it does not establish a failure of cubic metric rigidity in dimension
four.

\section{Prior results and the additional deduction}
The lower-dimensional comparison is prior work: cubic surfaces by
Odaka--Spotti--Sun, cubic threefolds by Liu--Xu, and cubic fourfolds by
Liu \cite{OSS,LX,Liu}. Cubic curves are not included in a positive-KE
Fano headline. Smooth KE existence should not be advertised as newly
discovered here.

The primary author page \cite{KSZZ} currently links a September 30, 2026
manuscript, publicly downloaded on October 6, 2026 (Pacific time).
Its Theorem 1.1 states the arbitrary-dimensional cubic comparison,
including stacks and good moduli spaces; its proof uses refined local
inequalities and tangent-sheaf semistability. The actual earliest web
posting is unverified, and a manuscript date alone does not establish
public priority. Its present public availability is verified. This
audit does not independently certify every lemma of that separate paper.

The complex cubic comparison is therefore prior public work. Our
gap-based account instantiates the known Spotti--Sun reduction; the
unrestricted gap proof belongs to OpenAI, and global Reeb minimization
to Li--Liu.

Theorem~\ref{thm:index} and the singular metric-fiber argument are the
additional deductions recorded here. Smooth cubic tangent stability
already follows from Peternell--Wi\'sniewski
\cite[Corollary 1.5 and Theorem 2.3]{PW}. Smooth positive-KE rigidity up
to factorwise conjugation is explained by Spotti \cite[Theorem 1.1.1]{Spotti}.
Odaka--Spotti--Sun expressly give the ordinary metric quotient by
conjugation for their singular del Pezzo-surface compactifications
\cite[\S\S1--2]{OSS}. The present argument checks the singular
higher-dimensional boundary by the Cartier-index exclusion of product
covers, reflexive endomorphisms, and the GH regular-locus theorem.
It uses DGP's splitting, not a newly invented decomposition principle.
The threshold's product counterexample is elementary.

Primary-source comparisons distinguish these exact scopes; they do
not certify first public disclosure of every possible equivalent
corollary. We make no first-priority claim and do not advertise a new
solution of the complex cubic K/GIT problem. The contribution is a
modest, explicitly attributed singular metric-rigidity criterion and
its boundary application.

\begin{thebibliography}{99}
\bibitem{OAI} OpenAI, \emph{The ordinary-double-point gap in every dimension},
OpenAI Math Release preprint, September 24, 2026.
\href{https://github.com/openai/math/blob/adc7f1241b42e322a6451854ab7e4b4c146bf78a/preprints/The-ordinary-double-point-gap-in-every-dimension-September-24-2026/paper.pdf}{Pinned manuscript}.
Sole companion: \emph{The normalized-volume gap in dimension four}.
\bibitem{LL} C. Li and Y. Liu, \emph{K\"ahler--Einstein metrics and volume
minimization}, Adv. Math. 341 (2019), 440--492,
\href{https://arxiv.org/abs/1602.05094}{arXiv:1602.05094v3}.
\bibitem{SS} C. Spotti and S. Sun, \emph{Explicit Gromov--Hausdorff
compactifications of moduli spaces of K\"ahler--Einstein Fano manifolds},
Pure Appl. Math. Q. 13 (2017), 477--515,
\href{https://arxiv.org/abs/1705.00377}{arXiv:1705.00377v1}.
\bibitem{vC} C. van Coevering, \emph{Ricci-flat K\"ahler metrics on
crepant resolutions of K\"ahler cones},
\href{https://arxiv.org/abs/0806.3728}{arXiv:0806.3728v3}.
\bibitem{CS} T. C. Collins and G. Sz\'ekelyhidi, \emph{Sasaki--Einstein
metrics and K-stability}, Geom. Topol. 23 (2019), 1339--1413,
\href{https://doi.org/10.2140/gt.2019.23.1339}{doi:10.2140/gt.2019.23.1339}.
\bibitem{Fujita} T. Fujita, \emph{On singular del Pezzo varieties}, in
\emph{Algebraic Geometry}, Lecture Notes in Math. 1417,
Springer, 1990, 117--128,
\href{https://doi.org/10.1007/BFb0083337}{doi:10.1007/BFb0083337}.
\bibitem{LWX} C. Li, X. Wang and C. Xu, \emph{On the proper moduli spaces
of smoothable K\"ahler--Einstein Fano varieties}, Duke Math. J. 168 (2019),
1387--1459, \href{https://arxiv.org/abs/1411.0761}{arXiv:1411.0761v4}.
\bibitem{OSS} Y. Odaka, C. Spotti and S. Sun, \emph{Compact moduli spaces
of del Pezzo surfaces and K\"ahler--Einstein metrics}, J. Differential Geom.
102 (2016), 127--172, \href{https://arxiv.org/abs/1210.0858}{arXiv:1210.0858}.
\bibitem{LX} Y. Liu and C. Xu, \emph{K-stability of cubic threefolds},
Duke Math. J. 168 (2019), 2029--2073,
\href{https://arxiv.org/abs/1706.01933}{arXiv:1706.01933}.
\bibitem{Liu} Y. Liu, \emph{K-stability of cubic fourfolds},
J. Reine Angew. Math. 786 (2022), 55--77,
\href{https://arxiv.org/abs/2007.14320}{arXiv:2007.14320v2}.
\bibitem{KSZZ} B. Kong, Y. Shen, J. Zhao and M. Zhao,
\emph{K-moduli of cubic hypersurfaces}, manuscript dated September 30, 2026,
\href{https://sites.google.com/uic.edu/jzhao/research}{primary author page},
\href{https://drive.google.com/file/d/1eB16wLGE2G5QrbV-pUDmw_ieLMiLlOcO/view}{public PDF}.
Access checked October 6, 2026 Pacific time.
\bibitem{DGP} S. Druel, H. Guenancia and M. P\u aun,
\emph{A decomposition theorem for $\Q$-Fano K\"ahler--Einstein varieties},
Comptes Rendus Math. 362, Special issue S1 (2024), 93--118,
\href{https://doi.org/10.5802/crmath.612}{doi:10.5802/crmath.612}.
\bibitem{DSII} S. Donaldson and S. Sun,
\emph{Gromov--Hausdorff limits of K\"ahler manifolds and algebraic
geometry, II}, J. Differential Geom. 107 (2017), 327--371,
\href{https://arxiv.org/abs/1507.05082}{arXiv:1507.05082v1}.
\bibitem{SGA2} A. Grothendieck,
\emph{Cohomologie locale des faisceaux coh\'erents et th\'eor\`emes de
Lefschetz locaux et globaux (SGA 2)}, North-Holland, 1968;
corrected edition by M. Raynaud,
\href{https://arxiv.org/abs/math/0511279}{arXiv:math/0511279}.
\bibitem{Stacks} The Stacks Project Authors, \emph{The Stacks Project},
\href{https://stacks.math.columbia.edu/tag/0BEM}{tag 0BEM},
Lemma 33.45.1; \href{https://stacks.math.columbia.edu/tag/0BEL}{tag 0BEL},
Section 33.45. Accessed October 6, 2026 Pacific time.
\bibitem{PW} T. Peternell and J. A. Wi\'sniewski,
\emph{On stability of tangent bundles of Fano manifolds with $b_2=1$},
J. Algebraic Geom. 4 (1995), 363--384,
\href{https://arxiv.org/abs/alg-geom/9306010}{arXiv:alg-geom/9306010v2}.
\bibitem{Spotti} C. Spotti, \emph{Degenerations of K\"ahler--Einstein
Fano manifolds}, doctoral thesis, 2012,
\href{https://arxiv.org/abs/1211.5334}{arXiv:1211.5334}.
\bibitem{KV} C. Hacon, \emph{Math 7800 notes}, 2018,
\href{https://www.math.utah.edu/~hacon/7800/Math7800-2018.pdf}{primary notes},
Theorem 2.22 (klt Kawamata--Viehweg vanishing).
\bibitem{GAGA} J.-P. Serre, \emph{G\'eom\'etrie alg\'ebrique et
g\'eom\'etrie analytique}, Ann. Inst. Fourier 6 (1956), 1--42,
\href{https://www.numdam.org/article/AIF_1956__6__1_0.pdf}{primary article}.
\end{thebibliography}
\end{document}
